\documentclass[11pt]{amsart}

\usepackage[margin=1in]{geometry}
\usepackage{amsmath,amssymb,amsthm,mathtools}
\usepackage{enumitem}
\usepackage[hidelinks]{hyperref}
\usepackage{microtype}

\setlength{\parskip}{0.6em}
\setlength{\parindent}{0pt}

\allowdisplaybreaks
\setlist[itemize]{topsep=2pt,itemsep=1pt,leftmargin=1.5em}
\setlist[enumerate]{topsep=2pt,itemsep=1pt,leftmargin=1.6em}

\newtheorem{theorem}{Theorem}[section]
\newtheorem{proposition}[theorem]{Proposition}
\newtheorem{lemma}[theorem]{Lemma}
\newtheorem{corollary}[theorem]{Corollary}
\theoremstyle{definition}
\newtheorem{definition}[theorem]{Definition}
\theoremstyle{remark}
\newtheorem{remark}[theorem]{Remark}

\newcommand{\nuu}{\nu_2}
\newcommand{\Ecal}{\mathcal E}
\newcommand{\Rcal}{\mathcal R}
\newcommand{\V}{\mathcal V}
\newcommand{\GD}{G_D}
\newcommand{\CD}{C_D}
\newcommand{\bdD}{\partial_D}
\newcommand{\depth}{\operatorname{depth}}
\newcommand{\dist}{\operatorname{dist}}
\newcommand{\oddrep}{\operatorname{oddrep}}
\newcommand{\Pre}{\mathsf{Pre}}
\newcommand{\Seed}{\operatorname{Seed}}
\newcommand{\Pair}{\mathsf P}
\newcommand{\Sfam}{\mathcal S}
\newcommand{\eps}{\varepsilon}

\title{A finite graph from the 2-adic valuation of the \texorpdfstring{$3n+1$}{3n+1} map}
\author{Aditya Sriram}

\address{Independent Researcher}
\email{adi.sriram.math@gmail.com}
\date{April 2026}

\begin{document}
\begin{abstract}
\setlength{\parskip}{0.6em} 
\setlength{\parindent}{0pt} 

This paper introduces a finite directed graph $G_D$ inspired by the $3n+1$ map.

The vertices are pairs $(r,d)$, where $1\le d\le D$ and $r$ is an odd 
residue modulo $2^d$. An edge $(r,d) \longrightarrow (s,e)$ encodes the
transition to the next odd residue class $(s,e)$ via the valuation data 
at $(r,d)$. Each such edge carries a weight that measures the logarithmic 
drift of the odd-to-odd map.

The main result is that $G_D$ has a unique nontrivial strongly connected
component $C_D$ of size $D(D-2)$ for $D\ge 3$, and each depth-$d$,
distance-$k$ layer outside $C_D$ is in bijection with 
\[
\Sfam_{d,k} = \{\Sigma\subseteq[d] : d\in\Sigma,\ |\Sigma|=2k+2\}.
\]
\end{abstract}
\maketitle


\section{Introduction}\label{sec:introduction}
The $3n+1$ problem is easy to state.  Given a positive integer $n$, define
\[
T(n)=
\begin{cases}
(3n+1)/2, & n\text{ odd},\\
n/2, & n\text{ even}.
\end{cases}
\]
Pick any positive integer.  Iterate.  The conjecture is that you always
reach~$1$.  Nobody has proved it.  The standard references are Lagarias'
survey~\cite{Lagarias1985} and the annotated bibliographies that
followed~\cite{LagariasAnnotated1999, LagariasAnnotated2009, Lagarias2010}.

One natural picture comes from running the map backwards.  Every integer has a
unique successor, so the predecessors form a tree.  Its branching and
stopping-time statistics have been studied
extensively~\cite{ApplegateLagarias1995, ApplegateLagarias2003, Wirsching1998}.

A second picture comes from passing to the $2$-adic integers.  The odd-to-odd
map
\[
S(x)=\frac{3x+1}{2^{\nu_2(3x+1)}}
\]
extends continuously to $\mathbb{Z}_2$.  Bernstein and Lagarias constructed a
conjugacy map that intertwines the $2$-adic shift with $S$ and analyzed the
induced permutations modulo~$2^n$~\cite{Bernstein1994,
BernsteinLagarias1996}.  The $2$-adic valuation $\nu_2(3x+1)$ is the
structural datum throughout this line of work.

Now here is the gap.

In integer space, the predecessor picture is a tree.  In $2$-adic residue
space, it isn't.  A residue class modulo~$2^d$ can have multiple predecessors
at higher depth, and the dynamics between classes can cycle.  Laarhoven and
de~Weger made this precise: at a fixed modulus~$2^d$, the $3n+1$ graph on
residue classes is isomorphic to the binary de~Bruijn graph, and the infinite
$2$-adic limit recovers the Bernstein--Lagarias
conjugacy~\cite{LaarhovenDeWeger2013}.

But that fixes the modulus.  What happens when a residue class that cannot yet
determine its transition is refined to the next depth instead?  I did not find
that question taken up in the literature.

So I ran the experiment.  At depth~$d$, let a vertex $(r,d)$ represent all odd
integers congruent to $r$ modulo~$2^d$.  Sometimes the class already
determines the next odd-to-odd transition completely, and it gets a forward
edge to a lower-depth class.  Sometimes the visible bits are not enough, and
the class must be split at depth $d+1$.  The result is a finite directed graph
$G_D$ on vertices at varying depths.  It is neither the predecessor tree nor
the fixed-modulus de~Bruijn graph.  It has cycles.

The structure turned out to be rigid.  At each depth~$d$, exactly $d$ residue
classes are \emph{exceptional}: the ones where the valuation data are not yet
resolved.  That count being exactly~$d$ was not assumed; it emerged from initial
computation and is proved in Section~\ref{sec:core}.

The exceptional residues have an explicit algebraic form. Their
children at depth $d+1$ split one-for-one into a new exceptional residue and a
regular residue that maps straight to the base class $(1,1)$. The recurrent
part of the graph assembles into a single strongly connected component of
size~$D(D-2)$.

The valuation data turned out to be load-bearing.  The pair of exponents that
determine each edge also carry the Syracuse log-drift that appears throughout
the $3n+1$ literature~\cite{Terras1976, Everett1977, Eliahou1993, Tao2022};
the precise connection is worked out in Section~\ref{sec:collatz}.

Representing an infinite-set problem in a finite structure buys you nothing
unless the finite object comes with a natural correspondence to the data you
want to track~\cite{Stanley2012, FlajoletSedgewick2009}.  By proving the
bijection, structural facts about the adaptive two-adic residue dynamics become
directly readable from the combinatorics of subsets of~$[d]$.  For each
cutoff~$D$, the graph~$G_D$ has a unique nontrivial strongly connected
component~$C_D$, and every vertex outside~$C_D$ carries a canonical coordinate
built from iterated valuation-preimage pairs.  After the natural top-bit
quotient, the depth-$d$, distance-$k$ layer is in canonical bijection with
\[
\Sfam_{d,k}=\{\Sigma\subseteq[d]:d\in\Sigma,\ |\Sigma|=2k+2\}.
\]
$3n+1$ is the origin of the construction, but what is proved is a finite
graph-theoretic and enumerative statement about adaptive two-adic valuation
data.  It identifies a finite object, classifies it completely, and shows that
its weighted structure faithfully encodes the logarithmic drift that runs
through the $3n+1$ literature.

\section{The valuation signature}\label{sec:syracuse}
The graph encodes one arithmetic operation. The grouped odd-to-odd
Collatz step and its two-adic bookkeeping. Both ideas are standard,
see Wirsching~\cite{Wirsching1998}.

Given an odd integer $r$, write $r+1=2^tq$ with $q$ odd.  The grouped
step produces
\[
3^t q - 1 = 2^\rho Y,
\qquad
\text{where } Y \text{ is odd and } \rho = \nu_2(3^t q - 1).
\]

We call the pair $(t,\rho)$ the \emph{valuation signature}: $t$
halvings from the binary structure of $r+1$, and $\rho$ more stripped
after the multiplication.

Now suppose $r$ is known only modulo $2^d$.  You have $d$ bits of
precision to work with.  The class pins down $t$ exactly when
$\nu_2(r+1)<d$, and pins down $\rho$ exactly when
$3^tq\not\equiv1\pmod{2^{d-t}}$.  When both are determined, the
target depth is
\[
e=d-t-\rho.
\]
The $d$ bits of input precision are consumed: $t$ by the structure of
$r+1$ and $\rho$ by the structure of $3^tq-1$.  What remains is $e$.

The graph is built from this quantity.

\section{The graph}\label{sec:graph}
\subsection{Vertices}
A vertex of $G_D$ is a pair $(r,d)$ with $1\le d\le D$ and $r$ an odd
residue modulo $2^d$.  When $e<1$, the class has not resolved the
valuation signature. It is replaced by:
\[
(r,d)\;\longmapsto\;(r,d+1),\quad (r+2^d,d+1).
\]

\subsection{Edges}
When $e\ge1$, the grouped step is well-defined.  Set
$y=(3^tq-1)/2^\rho$ and $s=\oddrep_e(y)$, the odd representative
of $y$ modulo $2^e$.  The edge is
\[
(r,d)\to(s,e).
\]

\section{Local structure}
This section proves the structural constraint behind the bijection: unresolved residue classes split predictably and subsequently resolve.

The condition $e\le0$ means $3^tq\equiv1\pmod{2^{d-t}}$: the
remaining $d-t$ bits of precision are exactly absorbed, and there is
nothing left to read $\rho$ from.  Since $r+1=2^tq$, this forces
\[
r_{d,t}\equiv 2^t(3^t)^{-1}-1\pmod{2^d},
\qquad 1\le t<d.
\]
Set $r_{d,d}=2^d-1$.  These are the
\emph{unresolved residues}
$\Ecal_d=\{r_{d,1},\ldots,r_{d,d}\}$, and they are distinct since
$\nuu(r_{d,t}+1)=t$.  Every other residue is \emph{resolved}.

An unresolved residue of type $t<d$ is \emph{non-maximal}; the
remaining case $r_{d,d}=2^d-1$ is \emph{maximal}.

\begin{lemma}[Direct returns]\label{lem:directreturns}
Every non-maximal unresolved vertex, when refined, produces one child
that stays unresolved and one resolved child that maps to $(1,1)$.
The maximal vertex produces two unresolved children.
Call the resolved child a \emph{direct return}.  At depth $d$ there are
exactly $\max(d-2,0)$ direct returns.
\end{lemma}

\begin{proof}
For the non-maximal case, the two depth-$(d+1)$ lifts of
$(r_{d,t},d)$ replace $q$ by $q$ and $q+2^{d-t}$.  Since
\[
3^t(q+2^{d-t})-1=(3^tq-1)+3^t\cdot 2^{d-t},
\]
exactly one satisfies the stronger congruence modulo $2^{d+1-t}$;
that is the unresolved child of type~$t$.  The other has $\rho=d-t$,
giving target depth~$1$.

For the maximal case, the two children of $(2^d-1,d)$ are $2^d-1$ and
$2^{d+1}-1$: the unresolved residues of types $d$ and $d+1$ at
depth $d+1$.  Neither is a direct return.

The non-maximal unresolved vertices at depth $d-1$ are those of types
$1,\ldots,d-2$, each contributing one direct return at depth~$d$.
The count is $\max(d-2,0)$.
\end{proof}

\section{The unique nontrivial SCC}\label{sec:core}
The graph $G_D$ has one nontrivial strongly connected component.
Every vertex either belongs to it, flows into it, or is trapped on
the boundary.  This section identifies the component, counts it, and
shows that no other cycle exists.

First, the boundary.  Lemma~\ref{lem:directreturns} guarantees that
every unresolved vertex produces at least one unresolved child at the
next depth.  At depth $D$, there is no next depth.  The unresolved
vertices of types $1,\ldots,D-1$ have no outgoing valuation edges and
no room to refine; call them \emph{terminal}.  The vertex
$(2^{D-1}-1,D-1)$ refines only to terminal vertices.  Together these
$D$ vertices form the boundary $\bdD$.  Nothing leaves it.

\begin{theorem}[SCC theorem]\label{thm:core}
For $D\ge3$, $G_D$ has a unique nontrivial strongly connected
component $C_D$.  It consists of all unresolved vertices and direct
returns outside the boundary, and
\[
|C_D|=D(D-2).
\]
A vertex reaches $C_D$ if and only if it is not in $\bdD$.
\end{theorem}

\begin{proof}
The argument has three parts: the claimed set is strongly connected,
everything outside it flows in, and no other cycle exists.

\medskip\noindent\emph{Strong connectivity.}
Every direct return maps to $(1,1)$.  Every non-maximal unresolved
vertex in $C_D$ has a direct-return child
(Lemma~\ref{lem:directreturns}), and hence a path to $(1,1)$.  Every
maximal unresolved vertex in $C_D$ has depth at most $D-2$, since it
is not on the boundary; one refinement step produces a non-maximal
unresolved child, and the previous case applies.  Conversely, $(1,1)$
reaches every unresolved vertex by repeatedly taking unresolved parents
--- the defining congruence at depth $d$ restricts to the same
congruence at depth $d-1$ --- and every direct return is a refinement
child of an unresolved vertex.  So $(1,1)$ reaches all of $C_D$.

\medskip\noindent\emph{Absorption.}
Let $v$ be a resolved vertex that is not a direct return.  It has a
valuation edge: the condition $e\le0$ would force $v$ to be unresolved.
That edge strictly lowers depth, since $t\ge1$ and $\rho\ge1$.
Iterating reaches either $C_D$ or the boundary.  It cannot hit the
boundary first: every valuation edge from depth at most $D$ has target
depth at most $D-2$, while the boundary sits at depths $D-1$ and $D$.

\medskip\noindent\emph{No other cycles.}
Outside $C_D$, every path either enters the boundary or follows
valuation edges of strictly decreasing depth.  Neither admits a cycle.

\medskip\noindent\emph{Count.}
The unresolved vertices across all depths number
$\sum_{d=1}^D d = D(D+1)/2$.  The direct returns number
$\sum_{d=3}^D(d-2) = (D-1)(D-2)/2$.  Their union has size $D^2-D+1$.
The boundary removes $D$ terminal unresolved vertices and one
additional vertex, giving $|C_D|=D^2-D+1-(D+1)=D(D-2)$.
\end{proof}

\section{Preimages and pair indexing}\label{sec:preimages}
Fix a target vertex and solve for all valuation-edge sources.  A
preimage is not indexed by an arbitrary residue.  It is indexed by the
two new depth coordinates it inserts between the target depth and the
source depth.

\begin{theorem}[Preimage-pair theorem]\label{thm:preimage}
Fix a target vertex $u=(s,e)\in V(G_D)$.  The valuation-edge preimages
of $u$ are canonically indexed by pairs
\[
e<i<j\le D.
\]
For the pair $(i,j)$, set
\[
t=i-e,
\qquad
\rho=j-i,
\qquad
 d=j.
\]
Then the unique source is obtained by solving
\[
q\equiv3^{-t}(1+2^\rho s)\pmod {2^{e+\rho}}
\]
and setting
\[
\Pre_{i,j}(s,e)=\bigl(\oddrep_j(2^tq-1),j\bigr).
\]
Consequently $|F^{-1}(s,e)|=\binom{D-e}{2}$.
\end{theorem}

\begin{proof}
If a valuation edge has source depth $d$ and target depth $e$, then
\[
d=e+t+\rho.
\]
Thus $i=e+t$ and $j=d$ satisfy $e<i<j\le D$.

Conversely, fix $e<i<j\le D$ and put $t=i-e$, $\rho=j-i$.  A source
with target residue $s$ must satisfy
\[
3^tq-1\equiv2^\rho s\pmod {2^{e+\rho}},
\]
or equivalently
\[
q\equiv3^{-t}(1+2^\rho s)\pmod {2^{e+\rho}}.
\]
This congruence has a unique odd solution modulo $2^{e+\rho}$.  With
$r=\oddrep_j(2^tq-1)$, one has $\nuu(r+1)=t$.  Since $s$ is odd, the
valuation of $3^tq-1$ is exactly $\rho$, and division by $2^\rho$
gives target residue $s$ modulo $2^e$.  The source is resolved because
$e\ge1$.  Uniqueness follows from the same congruence for $q$.
\end{proof}

The pair attached to a valuation edge $v=(r,d)\to(s,e)$ is therefore
\[
\Pair(v)=\{e+\nuu(r+1),d\}.
\]
Each preimage step adds two depths: an intermediate valuation depth and
the source depth.

\section{Histories and binomial layers}\label{sec:histories}
Let
\[
R_D(v)=\dist(v,C_D)
\]
when $v$ reaches $C_D$, and $R_D(v)=\infty$ otherwise.  If
$R_D(v)=k>0$, then $v$ is resolved and not a direct return, and
repeated valuation edges give a unique chain
\[
v=v_k\to v_{k-1}\to\cdots\to v_0\in C_D.
\]


For $d\ge2$, let $\tau_d$ be the top-bit involution
\[
\tau_d(r,d)=
\begin{cases}
(r+2^{d-1},d),&r<2^{d-1},\\
(r-2^{d-1},d),&r>2^{d-1}.
\end{cases}
\]
It exchanges the two depth-$d$ lifts of the same residue modulo $2^{d-1}$.

\begin{proposition}[Stable core seeds]\label{prop:seeds}
For $2\le d\le D-2$, the quotient $(C_D\cap V_d)/\tau_d$ is canonically labeled by the two-element sets
\[
\{a,d\},\qquad 1\le a<d.
\]
There are two oriented core vertices above each label.
\end{proposition}

\begin{proof}
For $1\le a\le d-2$, the two children of the exceptional vertex of type $a$ at depth $d-1$ consist of the exceptional vertex of type $a$ at depth $d$ and its direct-return sibling.  These two children differ by $2^{d-1}$, so they form one $\tau_d$-orbit.  The remaining orbit comes from the two children of the maximal exceptional vertex at depth $d-1$, namely the exceptional vertices of types $d-1$ and $d$.  Since $d\le D-2$, no boundary vertex is involved.  These $d-1$ two-element orbits exhaust $C_D\cap V_d$.
\end{proof}

Here \emph{stable} means that the depth is away from the truncation boundary: depths $2\le d\le D-2$ have the same core-seed structure they would have in any larger cutoff graph.

\begin{definition}[Seed and orientation]\label{def:orientation}
For a stable core vertex $v_0=(r,d)\in C_D$ with $2\le d\le D-2$, let
\[
\Seed(v_0)=\{a,d\}
\]
be the label of its $\tau_d$-orbit from Proposition~\ref{prop:seeds}.  Its orientation is the top-bit sign
\[
\eps(v_0)=
\begin{cases}
-,&1\le r<2^{d-1},\\
+,&2^{d-1}<r<2^d.
\end{cases}
\]
For any positive-distance vertex $v$, define $\eps(v)$ to be the orientation of the terminal core vertex $v_0$ in its valuation chain.
\end{definition}

If $v=v_k\to\cdots\to v_0\in C_D$ is the unique chain for a vertex of positive finite distance, write $v_m=(r_m,d_m)$.  The endpoint $v_0$ has stable depth: $2\le d_0\le D-2$, because the final edge is not a direct return and every valuation edge lowers depth by at least two.  Define the projected history
\[
\Sigma(v)=\Seed(v_0)\cup\bigcup_{m=1}^k\{d_{m-1}+\nuu(r_m+1),d_m\}.
\]
For a stable core vertex itself, set $\Sigma(v_0)=\Seed(v_0)$.

\begin{lemma}[History intervals]\label{lem:historyintervals}
If $v$ has finite distance $k$ and depth $d$, then the union defining $\Sigma(v)$ is disjoint.  For $k\ge1$,
\[
\Sigma(v)\subseteq[d],\qquad d\in\Sigma(v),\qquad |\Sigma(v)|=2k+2.
\]
\end{lemma}

\begin{proof}
For the edge $v_m\to v_{m-1}$, put
\[
t_m=\nuu(r_m+1),\qquad \rho_m=\nuu(3^{t_m}q_m-1).
\]
Then
\[
d_m=d_{m-1}+t_m+\rho_m,
\]
so
\[
\{d_{m-1}+t_m,d_m\}\subset(d_{m-1},d_m].
\]
The seed lies in $[1,d_0]$.  Inducting along the chain, each preimage step appends two new elements in the next interval $(d_{m-1},d_m]$.  Hence the union is disjoint, lies in $[d]$, contains the final depth $d$, and has cardinality $2+2k$.
\end{proof}

For $k\ge1$, define
\[
\Sfam_{d,k}=\{\Sigma\subseteq[d]:d\in\Sigma,\ |\Sigma|=2k+2\}.
\]

\begin{theorem}[History bijection]\label{thm:bijection}
For $D\ge3$, $k\ge1$, and $1\le d\le D$, the map
\[
v\longmapsto (\eps(v),\Sigma(v))
\]
is a canonical bijection
\[
\{v:\depth(v)=d,\ R_D(v)=k\}
\longrightarrow
\{+,-\}\times\Sfam_{d,k}.
\]
Here canonical means that the inverse is computed by the explicit preimage construction of Theorem~\ref{thm:preimage}.  Consequently
\[
\bigl|\{v:\depth(v)=d,\ R_D(v)=k\}\bigr|=2\binom{d-1}{2k+1}.
\]
\end{theorem}

\begin{proof}
Lemma~\ref{lem:historyintervals} shows that the map lands in $\{+,-\}\times\Sfam_{d,k}$.

To reconstruct, take $(\eps,\Sigma)$ with
\[
\Sigma=\{\sigma_1<\sigma_2<\cdots<\sigma_{2k+2}=d\}.
\]
Start with the oriented core vertex labeled by $\{\sigma_1,\sigma_2\}$ and orientation $\eps$.  At step $m=1,\ldots,k$, use Theorem~\ref{thm:preimage} to take the unique valuation-edge preimage of the current target of depth $\sigma_{2m}$ indexed by the pair
\[
\sigma_{2m}<\sigma_{2m+1}<\sigma_{2m+2}.
\]
This constructs a chain whose final source has depth $d$.

Each reconstructed edge has target depth at least $2$, so no reconstructed source is a direct return; hence the first time the chain reaches the core is its prescribed endpoint.  Thus the final source has distance $k$.  The construction is inverse to the history map because Theorem~\ref{thm:preimage} gives a unique source for each target and pair.

The count follows from
\[
|\Sfam_{d,k}|=\binom{d-1}{2k+1},
\]
since $d$ is forced to lie in $\Sigma$.
\end{proof}

\begin{lemma}[Top-bit equivariance]\label{lem:topbit}
Let $v=(r,d)$ be a regular vertex with valuation edge $F(v)=(s,e)$ and suppose $e\ge2$.  Then $\tau_d(v)$ is regular with the same valuation labels $(t,\rho)$, and
\[
F(\tau_d v)=\tau_e F(v).
\]
\end{lemma}

\begin{proof}
Write
\[
t=\nuu(r+1),\qquad q=\frac{r+1}{2^t},\qquad \rho=\nuu(3^tq-1),
\]
so $d=e+t+\rho$.  Replacing $r$ by its top-bit flip changes $r+1$ to
\[
r+1\pm2^{d-1}.
\]
Since $e\ge2$ and $\rho\ge1$, one has $d-1>t$, so the value of $t$ is unchanged.  The corresponding odd quotient is
\[
q'=q\pm2^{d-1-t}.
\]
Then
\[
3^tq'-1=(3^tq-1)\pm3^t2^{d-1-t}.
\]
The added term has two-adic valuation $d-1-t=\rho+e-1>\rho$, so the exact value of $\rho$ is also unchanged.  After division by $2^\rho$, the target odd representative changes by
\[
\pm 3^t2^{e-1},
\]
which is congruent modulo $2^e$ to adding or subtracting the top bit $2^{e-1}$.  Hence the target is $\tau_e(s,e)$.
\end{proof}

\begin{corollary}[Quotient layers]\label{cor:quotient}
For $k\ge1$ and $d\ge2$,
\[
\{v:\depth(v)=d,\ R_D(v)=k\}/\tau_d
\cong
\Sfam_{d,k}.
\]
Equivalently, summing over $d\le D$,
\[
\bigl|\{v:R_D(v)=k\}\bigr|=2\binom{D}{2k+2}.
\]
\end{corollary}

\begin{proof}
Positive-distance valuation chains never use a direct-return edge, and every target depth along such a chain is at least $2$.  Lemma~\ref{lem:topbit} therefore shows that the whole chain commutes with the top-bit involution.  The involution changes only the orientation coordinate $\eps$ and leaves the subset history $\Sigma$ unchanged.  Quotienting Theorem~\ref{thm:bijection} by this two-element orientation gives the first claim.

For the total count, sum the oriented layer count:
\[
\sum_{d=1}^D2\binom{d-1}{2k+1}
=
2\sum_{d=2k+2}^D\binom{d-1}{2k+1}
=
2\binom{D}{2k+2},
\]
where the last equality is the hockey-stick identity.
\end{proof}

\section{Weighted histories}\label{sec:weighted}

The history coordinate is not only a counting device.  It also records the additive statistic that corresponds to Syracuse logarithmic drift.  In this section we regard $G_D$ as a weighted directed graph by assigning a real weight to each valuation edge from its valuation label.

Put
\[
\alpha=\log_2 3.
\]
If $v=(r,d)$ is regular and $F(v)=u=(s,e)$, write
\[
t=\nuu(r+1),\qquad q=\frac{r+1}{2^t},\qquad \rho=\nuu(3^tq-1).
\]
Thus $d=e+t+\rho$.  Define the edge weight
\[
\omega(v,u)=t\alpha-(t+\rho)=t\alpha-(d-e).
\]
Equivalently, if the pair of the edge is
\[
\Pair(v)=\{i,j\},\qquad e<i<j=d,
\]
then
\[
\omega(v,u)=(i-e)\alpha-(j-e).
\]

Now let $v$ have positive finite distance $R_D(v)=k$, and write its valuation chain as
\[
v=v_k\to v_{k-1}\to\cdots\to v_0\in C_D.
\]
The total chain weight is
\[
\Omega(v)=\sum_{m=1}^k \omega(v_m,v_{m-1}).
\]
The endpoint core vertex carries no weight in this definition; only the valuation edges in the distance chain are counted.

For a history
\[
\Sigma=\{\sigma_1<\sigma_2<\cdots<\sigma_{2k+2}=d\}\in\Sfam_{d,k},
\]
define the alternating-gap sum
\[
T(\Sigma)=\sum_{m=1}^k(\sigma_{2m+1}-\sigma_{2m})
\]
and the history weight
\[
\Omega(\Sigma)=T(\Sigma)\alpha-(d-\sigma_2).
\]

\begin{theorem}[Weight factorization through histories]\label{thm:weightfactor}
Let $D\ge3$, $k\ge1$, and let $v$ be a vertex with $R_D(v)=k$ and $\depth(v)=d$.  Then
\[
\Omega(v)=\Omega(\Sigma(v)).
\]
Consequently the total chain weight depends only on the quotient history $\Sigma(v)$ and is independent of the orientation $\eps(v)$.
\end{theorem}

\begin{proof}
Write the distance chain as
\[
v=v_k\to v_{k-1}\to\cdots\to v_0\in C_D,
\]
and put $v_m=(r_m,d_m)$.  For the edge $v_m\to v_{m-1}$, let
\[
t_m=\nuu(r_m+1),\qquad q_m=\frac{r_m+1}{2^{t_m}},\qquad \rho_m=\nuu(3^{t_m}q_m-1).
\]
By the definition of the history coordinate, if
\[
\Sigma(v)=\{\sigma_1<\sigma_2<\cdots<\sigma_{2k+2}=d\},
\]
then the $m$th valuation preimage step has
\[
d_{m-1}=\sigma_{2m},\qquad d_{m-1}+t_m=\sigma_{2m+1},\qquad d_m=\sigma_{2m+2}.
\]
Therefore
\[
t_m=\sigma_{2m+1}-\sigma_{2m},
\qquad
 t_m+\rho_m=d_m-d_{m-1}=\sigma_{2m+2}-\sigma_{2m}.
\]
Summing the edge weights gives
\begin{align*}
\Omega(v)
&=\sum_{m=1}^k\bigl(t_m\alpha-(t_m+\rho_m)\bigr)\\
&=\alpha\sum_{m=1}^k(\sigma_{2m+1}-\sigma_{2m})
  -\sum_{m=1}^k(\sigma_{2m+2}-\sigma_{2m})\\
&=T(\Sigma(v))\alpha-(\sigma_{2k+2}-\sigma_2)\\
&=T(\Sigma(v))\alpha-(d-\sigma_2).
\end{align*}
This is exactly $\Omega(\Sigma(v))$.  Since $\Sigma(v)$ is unchanged by switching the top-bit orientation, the total chain weight is independent of $\eps(v)$.
\end{proof}

\begin{corollary}[Weighted quotient layers]\label{cor:weightedquotient}
For $k\ge1$ and $d\ge2$, the quotient bijection
\[
\{v:\depth(v)=d,\ R_D(v)=k\}/\tau_d\cong\Sfam_{d,k}
\]
is weight-preserving when the left side is weighted by total chain weight and $\Sfam_{d,k}$ is weighted by $\Omega(\Sigma)$.
\end{corollary}

\begin{remark}
The statistic $T(\Sigma)$ selects the odd-numbered gaps after the core seed depth $\sigma_2$, while $d-\sigma_2$ is the total depth gained after the seed.  Thus the weight is an alternating-gap functional on the same subsets that index the binomial distance layers:
\[
\Omega(\Sigma)
=
\alpha\bigl[(\sigma_3-\sigma_2)+(\sigma_5-\sigma_4)+\cdots+(\sigma_{2k+1}-\sigma_{2k})\bigr]
-(d-\sigma_2).
\]
This is the structural point of the weighted theorem: once the core seed is fixed, the log-drift sees only the odd-indexed inserted gaps and the total post-seed depth gain.  No information from the orientation coordinate is needed.
\end{remark}

\section{Origin of the rule and meaning of the weight}\label{sec:collatz}

The rule used to define $G_D$ is extracted from a residue-class model of the odd Collatz/Syracuse dynamics.  In that model one studies the odd-to-odd valuation update
\[
X+1=2^tq,
\qquad
3^tq-1=2^\rho Y.
\]
This is not the usual one-step map $X\mapsto (3X+1)/2^{\nu_2(3X+1)}$.  Rather, it is a grouped valuation step: the first equation records a run-length $t$ determined by the binary position of $X+1$, and the second equation records the power of two removed after multiplying the odd quotient by $3^t$.

The adaptive graph keeps exactly the residue-class information needed to know this grouped step.  If a depth-$d$ class determines $t$ and determines $\rho$ before the visible precision is exhausted, it has a valuation edge.  If the visible precision is exhausted, the class is exceptional and must be refined.  This is why exceptional vertices refine and regular vertices move to lower depth.

The same equations explain the edge weight.  Since
\[
Y=\frac{3^tX+3^t-2^t}{2^{t+\rho}},
\]
one has
\[
\log_2(Y/X)
=
t\log_2 3-(t+\rho)
+
\log_2\!\left(1+\frac{1-(2/3)^t}{X}\right).
\]
For large lifts $X$ in a fixed residue class, the last term is $O(1/X)$.  Thus the graph weight
\[
\omega=t\log_2 3-(t+\rho)
\]
is precisely the residue-class part of the Syracuse multiplicative drift.  Along a finite chain, $\Omega(v)$ is the sum of these local log-drift terms.  Theorem~\ref{thm:weightfactor} says that this finite residue-level drift is a function only of the subset history $\Sigma$.

This connects the paper to the standard logarithmic bookkeeping used throughout Collatz analysis.  In cycle and stopping-time arguments, one compares the accumulated power of $3$ with the accumulated power of $2$; in logarithmic form this is a sum of terms of the shape
\[
(\text{number of }3\text{-multiplications})\log_2 3
-
(\text{number of }2\text{-divisions}).
\]
The present paper does not add a convergence claim to that tradition.  Its contribution is to show that, for the finite adaptive valuation graph, the corresponding log-functional has an exact combinatorial coordinate: the binomial subset history.

The statements proved here therefore have a deliberately limited scope.  A vertex $(r,d)$ is a residue class, not an integer.  The integer $1$ is represented by an infinite compatible tower of residue classes, and integer dynamics require lift data beyond any fixed $G_D$.  What the graph supplies is a finite, classifiable object: adaptive two-adic refinement produces a recurrent core, pair-indexed valuation preimages, a top-bit symmetry, binomial distance layers, and a history-level log-drift functional.

\begin{thebibliography}{99}

\bibitem{Allouche1979}
J.-P. Allouche.
Sur la conjecture de ``Syracuse--Kakutani--Collatz''.
\emph{S\'eminaire de Th\'eorie des Nombres de Bordeaux}, exp. no.~9, 1978--1979.

\bibitem{ApplegateLagarias1995}
D. Applegate and J. C. Lagarias.
The distribution of $3x+1$ trees.
\emph{Experimental Mathematics} 4 (1995), no.~3, 193--209.

\bibitem{ApplegateLagarias2003}
D. Applegate and J. C. Lagarias.
Lower bounds for the total stopping time of $3x+1$ iterates.
\emph{Mathematics of Computation} 72 (2003), no.~242, 1035--1049.

\bibitem{Bernstein1994}
D. J. Bernstein.
A non-iterative $2$-adic statement of the $3n+1$ conjecture.
\emph{Proceedings of the American Mathematical Society} 121 (1994), no.~2, 405--408.

\bibitem{BernsteinLagarias1996}
D. J. Bernstein and J. C. Lagarias.
The $3x+1$ conjugacy map.
\emph{Canadian Journal of Mathematics} 48 (1996), no.~6, 1154--1169.

\bibitem{Conway1972}
J. H. Conway.
Unpredictable iterations.
In \emph{Proceedings of the 1972 Number Theory Conference}, University of Colorado, Boulder, 1972, 49--52.

\bibitem{Diestel2017}
R. Diestel.
\emph{Graph Theory}.
Graduate Texts in Mathematics 173, 5th ed., Springer, 2017.

\bibitem{Eliahou1993}
S. Eliahou.
The $3x+1$ problem: new lower bounds on nontrivial cycle lengths.
\emph{Discrete Mathematics} 118 (1993), no.~1--3, 45--56.

\bibitem{Everett1977}
C. J. Everett.
Iteration of the number-theoretic function $f(2n)=n$, $f(2n+1)=3n+2$.
\emph{Advances in Mathematics} 25 (1977), no.~1, 42--45.

\bibitem{FlajoletSedgewick2009}
P. Flajolet and R. Sedgewick.
\emph{Analytic Combinatorics}.
Cambridge University Press, 2009.

\bibitem{Garner1981}
L. E. Garner.
On the Collatz $3n+1$ algorithm.
\emph{Proceedings of the American Mathematical Society} 82 (1981), no.~1, 19--22.

\bibitem{Gouvea1997}
F. Q. Gouv\^ea.
\emph{$p$-adic Numbers: An Introduction}.
Universitext, 2nd ed., Springer, 1997.

\bibitem{Hercher2023}
C. Hercher.
There are no Collatz $m$-cycles with $m \leq 91$.
\emph{Journal of Integer Sequences} 26 (2023), Article 23.3.5.

\bibitem{Korec1994}
I. Korec.
A density estimate for the $3x+1$ problem.
\emph{Mathematica Slovaca} 44 (1994), no.~1, 85--89.

\bibitem{KrasikovLagarias2003}
I. Krasikov and J. C. Lagarias.
Bounds for the $3x+1$ problem using difference inequalities.
\emph{Acta Arithmetica} 109 (2003), no.~3, 237--258.

\bibitem{Lagarias1985}
J. C. Lagarias.
The $3x+1$ problem and its generalizations.
\emph{The American Mathematical Monthly} 92 (1985), no.~1, 3--23.

\bibitem{LagariasAnnotated1999}
J. C. Lagarias.
The $3x+1$ problem: an annotated bibliography (1963--1999).
Preprint, arXiv:math/0309224, 2003.

\bibitem{LagariasAnnotated2009}
J. C. Lagarias.
The $3x+1$ problem: an annotated bibliography, II (2000--2009).
Preprint, arXiv:math/0608208, 2006.

\bibitem{Lagarias2010}
J. C. Lagarias, editor.
\emph{The Ultimate Challenge: The $3x+1$ Problem}.
American Mathematical Society, Providence, RI, 2010.

\bibitem{LaarhovenDeWeger2013}
T. Laarhoven and B. de Weger.
The Collatz conjecture and de Bruijn graphs.
\emph{Indagationes Mathematicae} 24 (2013), no.~4, 971--983.

\bibitem{MonksYazinski2004}
K. G. Monks and J. Yazinski.
The autoconjugacy of the $3x+1$ function.
\emph{Discrete Mathematics} 275 (2004), no.~1--3, 219--236.

\bibitem{Rozier2019}
O. Rozier.
Parity sequences of the $3x+1$ map on the $2$-adic integers and Euclidean embedding.
\emph{Integers} 19 (2019), Paper A8.

\bibitem{Stanley2012}
R. P. Stanley.
\emph{Enumerative Combinatorics, Volume 1}.
Cambridge Studies in Advanced Mathematics 49, 2nd ed., Cambridge University Press, 2012.

\bibitem{Tao2022}
T. Tao.
Almost all orbits of the Collatz map attain almost bounded values.
\emph{Forum of Mathematics, Pi} 10 (2022), e12.

\bibitem{Terras1976}
R. Terras.
A stopping time problem on the positive integers.
\emph{Acta Arithmetica} 30 (1976), 241--252.

\bibitem{Wirsching1998}
G. J. Wirsching.
\emph{The Dynamical System Generated by the $3n+1$ Function}.
Lecture Notes in Mathematics 1681, Springer, 1998.

\end{thebibliography}
\end{document}
